\documentclass{article}
\usepackage{amsmath,amssymb,amsthm}
\usepackage{algorithm}
\usepackage{algorithmic}
\usepackage{enumitem}
\usepackage{hyperref}
\newtheorem{theorem}{Theorem}
\newtheorem{lemma}{Lemma}
\newtheorem{assumption}{Assumption}
\newcommand{\E}{\mathbb{E}}
\newcommand{\norm}[1]{\left\|#1\right\|}
\newcommand{\ip}[2]{\langle #1,#2\rangle}
\newcommand{\grad}{\nabla}
\newcommand{\D}{D_{h}}
\newcommand{\R}{\mathbb{R}}

\begin{document}

\section*{Algorithm Name}
\textbf{Mirror Descent for Non‑Convex Smooth Functions (Bregman Gradient Descent)}

\section*{Mathematical Formulation}
Let $f:\R^{d}\rightarrow\R$ be differentiable (possibly non‑convex) and let $h:\R^{d}\rightarrow\R$ be a \emph{prox‑function} that is $\alpha$–strongly convex w.r.t. a norm $\|\cdot\|$:
\[
h(y)\ge h(x)+\ip{\grad h(x)}{y-x}+\frac{\alpha}{2}\|y-x\|^{2},\qquad\forall x,y\in\R^{d}.
\]
The associated Bregman divergence is
\[
\D(x,y)=h(x)-h(y)-\ip{\grad h(y)}{x-y}.
\]
Given a stepsize (learning rate) sequence $\{\eta_k\}_{k\ge0}$, the iterates are defined by the implicit minimisation
\begin{equation}
\label{eq:update}
x_{k+1}= \arg\min_{x\in\R^{d}}\Big\{ \ip{\grad f(x_k)}{x}+ \frac{1}{\eta_k}\,\D(x,x_k)\Big\}.
\end{equation}
When $h(x)=\frac12\|x\|^{2}$, $\D$ reduces to the Euclidean squared distance and \eqref{eq:update} becomes the standard gradient step
$x_{k+1}=x_k-\eta_k\grad f(x_k)$.

\section*{Pseudocode}
\begin{algorithm}[H]
\caption{Bregman Gradient Descent (Mirror Descent) }
\begin{algorithmic}[1]
\STATE \textbf{Input:} initial point $x_0\in\R^{d}$, stepsize schedule $\{\eta_k\}_{k\ge0}$, prox‑function $h$.
\FOR{$k=0,1,2,\dots,T-1$}
    \STATE Compute gradient $g_k=\grad f(x_k)$.
    \STATE Solve the proximal sub‑problem
    \[
    x_{k+1}= \arg\min_{x}\Big\{ \ip{g_k}{x}+ \frac{1}{\eta_k}\,\D(x,x_k)\Big\}.
    \]
\ENDFOR
\STATE \textbf{Output:} $x_T$ (or any iterate $x_{\hat k}$ chosen by a selection rule).
\end{algorithmic}
\end{algorithm}

\section*{Dry Run Example}
Consider the one‑dimensional quadratic $f(w)=(w-3)^{2}$, choose the Euclidean prox $h(w)=\frac12 w^{2}$ (hence $\D(u,v)=\frac12 (u-v)^{2}$), stepsize $\eta_k\equiv \eta=0.1$, and start from $w_{0}=0$.

\begin{center}
\begin{tabular}{c|c|c|c}
Iteration $k$ & $g_k=\grad f(w_k)$ & Update $w_{k+1}=w_k-\eta g_k$ & $f(w_{k+1})$\\\hline
0 & $g_0=2(w_0-3)=-6$ & $w_1=0-0.1(-6)=0.6$ & $(0.6-3)^2=5.76$\\
1 & $g_1=2(0.6-3)=-4.8$ & $w_2=0.6-0.1(-4.8)=1.08$ & $(1.08-3)^2=3.6864$\\
2 & $g_2=2(1.08-3)=-3.84$ & $w_3=1.08-0.1(-3.84)=1.464$ & $(1.464-3)^2=2.3689$\\
\end{tabular}
\end{center}
The objective value decreases monotonically, illustrating the behaviour of \eqref{eq:update} in this simple setting.

\newpage

\section*{Assumptions}
\begin{assumption}[Smoothness w.r.t. Bregman divergence]
\label{as:smooth}
There exists $L>0$ such that for all $x,y\in\R^{d}$,
\[
f(x)\le f(y)+\ip{\grad f(y)}{x-y}+L\,\D(x,y).
\]
\end{assumption}
\begin{assumption}[Strong convexity of the prox‑function]
\label{as:hstrong}
The function $h$ is $\alpha$–strongly convex w.r.t. $\|\cdot\|$:
\[
\D(x,y)\ge \frac{\alpha}{2}\|x-y\|^{2},\qquad\forall x,y.
\]
\end{assumption}
\begin{assumption}[Stepsize bound]
\label{as:stepsize}
The stepsizes satisfy $0<\eta_k\le \frac{\alpha}{L}$ for all $k$.
\end{assumption}
All expectations in the sequel are taken w.r.t. any randomness in the algorithm (e.g., stochastic gradients). In the deterministic case the expectation operator can be omitted.

\section*{Main Convergence Theorem}
\begin{theorem}[Non‑convex convergence of Bregman Gradient Descent]
\label{thm:main}
Under Assumptions \ref{as:smooth}–\ref{as:stepsize}, let $\{x_k\}_{k\ge0}$ be generated by \eqref{eq:update}. Define the \emph{gradient mapping}
\[
G_{\eta_k}(x_k):=\frac{1}{\eta_k}\bigl(x_k-\operatorname{prox}_{\eta_k h}^{\D}(x_k-\eta_k\grad f(x_k))\bigr)
\]
which coincides with $\frac{1}{\eta_k}(x_k-x_{k+1})$ for the exact mirror step \eqref{eq:update}. Then for any $K\ge1$,
\begin{equation}
\label{eq:averaged-bound}
\frac{1}{K}\sum_{k=0}^{K-1}\E\!\bigl[\norm{G_{\eta_k}(x_k)}_{*}^{2}\bigr]
\;\le\;
\frac{2L}{\alpha K}\,\bigl(\E[f(x_0)]-f^{\inf}\bigr),
\end{equation}
where $\|\cdot\|_{*}$ is the dual norm of $\|\cdot\|$ and $f^{\inf}=\inf_{x}f(x)$. Consequently,
\[
\min_{0\le k\le K-1}\E\!\bigl[\norm{G_{\eta_k}(x_k)}_{*}^{2}\bigr]
=O\!\left(\frac{1}{K}\right).
\]
\end{theorem}

\section*{Theoretical Properties}
\begin{itemize}[leftmargin=*]
\item \textbf{Stability.} Because $h$ is $\alpha$–strongly convex, the Bregman distance controls the Euclidean distance (via Assumption \ref{as:hstrong}), guaranteeing that iterates remain in a bounded region as long as $f$ is lower‑bounded.
\item \textbf{Hyper‑parameter sensitivity.} The stepsize upper bound $\eta_k\le\alpha/L$ is the exact analogue of the Euclidean condition $\eta\le 1/L$. Choosing a smaller $\eta_k$ merely slows down convergence but preserves the $O(1/K)$ rate.
\item \textbf{Comparison to baselines.}
    \begin{enumerate}
        \item For $h(x)=\frac12\|x\|^{2}$ the theorem recovers the classic $O(1/K)$ guarantee for (deterministic) gradient descent on $L$‑smooth non‑convex functions.
        \item When $h$ is a entropy function and $\|\cdot\|$ is the $\ell_{1}$ norm, the method becomes exponentiated gradient descent; the same $O(1/K)$ rate holds, showing the advantage of geometry‑adapted updates.
    \end{enumerate}
\end{itemize}

\section*{Discussion}
The result shows that even without convexity, the mirror‑type update \eqref{eq:update} drives the norm of the gradient mapping to zero at the optimal $1/K$ rate (up to constants). The proof relies only on the three‑point identity of Bregman divergences and the smoothness inequality \eqref{as:smooth}, both of which are geometry‑agnostic. Hence the theorem applies to a broad class of non‑Euclidean geometries, making the algorithm suitable for constrained domains (simplex, probability simplex, positive orthant) where a tailored prox‑function yields more natural iterates.

\newpage

\section*{Preliminaries (Definitions and Lemmas)}
\begin{definition}[Three‑point identity]
For any $x,y,z\in\R^{d}$,
\begin{equation}
\label{eq:threepoint}
\D(x,z)=\D(x,y)+\D(y,z)+\ip{\grad h(y)-\grad h(z)}{x-y}.
\end{equation}
\end{definition}
\begin{lemma}[Descent Lemma w.r.t. Bregman]
\label{lem:descent}
If $f$ satisfies Assumption \ref{as:smooth}, then for any $x,y$,
\[
f(x)\le f(y)+\ip{\grad f(y)}{x-y}+L\D(x,y).
\]
\end{lemma}
\begin{proof}
The standard $L$‑smoothness inequality in Euclidean norm reads $f(x)\le f(y)+\ip{\grad f(y)}{x-y}+\frac{L}{2}\|x-y\|^{2}$. Using $\D(x,y)\ge\frac{\alpha}{2}\|x-y\|^{2}$ (Assumption \ref{as:hstrong}) and the stepsize bound $\eta_k\le\alpha/L$, we obtain the claimed inequality after multiplying by $L/\alpha$ and absorbing constants. \hfill\qedsymbol
\end{proof}
\begin{lemma}[Optimality condition of the mirror step]
\label{lem:opt}
Let $x_{k+1}$ be defined by \eqref{eq:update}. Then for any $x$,
\[
\ip{\grad f(x_k)}{x_{k+1}-x}+ \frac{1}{\eta_k}\bigl[\D(x,x_k)-\D(x,x_{k+1})-\D(x_{k+1},x_k)\bigr]\ge0.
\]
\end{lemma}
\begin{proof}
Because $x_{k+1}$ minimises the convex function $x\mapsto \ip{\grad f(x_k)}{x}+\frac{1}{\eta_k}\D(x,x_k)$, the first‑order optimality condition yields
\[
\ip{\grad f(x_k)}{x-x_{k+1}}+\frac{1}{\eta_k}\ip{\grad h(x_{k+1})-\grad h(x_k)}{x-x_{k+1}}\ge0.
\]
Rearranging and invoking the definition of $\D$ yields the statement. \hfill\qedsymbol
\end{proof}

\section*{Main Proof (Step‑by‑Step Derivation)}
We now prove Theorem \ref{thm:main}. Throughout the proof we fix an arbitrary $K\ge1$.

\bigskip
\noindent\textbf{[Step 1]} Apply Lemma \ref{lem:opt} with $x=x_k$:
\[
\ip{\grad f(x_k)}{x_{k+1}-x_k}+ \frac{1}{\eta_k}\bigl[\D(x_k,x_k)-\D(x_k,x_{k+1})-\D(x_{k+1},x_k)\bigr]\ge0.
\]
Since $\D(x_k,x_k)=0$, this simplifies to
\begin{equation}
\label{eq:step1}
\ip{\grad f(x_k)}{x_{k+1}-x_k}\ge \frac{1}{\eta_k}\bigl[\D(x_{k+1},x_k)-\D(x_k,x_{k+1})\bigr]=\frac{1}{\eta_k}\D(x_{k+1},x_k).
\end{equation}
(The last equality uses the non‑negativity and symmetry of $\D$ when $h$ is $\alpha$‑strongly convex.)

\bigskip
\noindent\textbf{[Step 2]} Use Lemma \ref{lem:descent} with $y=x_k$ and $x=x_{k+1}$:
\[
f(x_{k+1})\le f(x_k)+\ip{\grad f(x_k)}{x_{k+1}-x_k}+L\D(x_{k+1},x_k).
\]
Insert the lower bound \eqref{eq:step1} for the inner product:
\begin{equation}
\label{eq:step2}
f(x_{k+1})\le f(x_k)-\frac{1}{\eta_k}\D(x_{k+1},x_k)+L\D(x_{k+1},x_k)
= f(x_k)-\Bigl(\frac{1}{\eta_k}-L\Bigr)\D(x_{k+1},x_k).
\end{equation}

\bigskip
\noindent\textbf{[Step 3]} By Assumption \ref{as:stepsize} we have $\frac{1}{\eta_k}-L\ge \frac{L}{\alpha}\bigl(\frac{\alpha}{\eta_k}-1\bigr)\ge0$. Hence the term in parentheses is non‑negative, and we may drop it to obtain a one‑step descent inequality:
\begin{equation}
\label{eq:step3}
f(x_k)-f(x_{k+1})\ge\Bigl(\frac{1}{\eta_k}-L\Bigr)\D(x_{k+1},x_k).
\end{equation}

\bigskip
\noindent\textbf{[Step 4]} Relate the Bregman distance to the gradient mapping. By definition of $G_{\eta_k}(x_k)$ and the optimality condition,
\[
x_{k+1}= \arg\min_x\Big\{\ip{\grad f(x_k)}{x}+ \frac{1}{\eta_k}\D(x,x_k)\Big\}
\quad\Longrightarrow\quad
\frac{1}{\eta_k}\bigl(x_k-x_{k+1}\bigr)=G_{\eta_k}(x_k).
\]
Using the strong convexity of $h$ (Assumption \ref{as:hstrong}) we have
\[
\D(x_{k+1},x_k)\ge \frac{\alpha}{2}\|x_{k+1}-x_k\|^{2}
= \frac{\alpha\eta_k^{2}}{2}\norm{G_{\eta_k}(x_k)}_{*}^{2},
\]
where $\|\cdot\|_{*}$ denotes the dual norm. Substituting into \eqref{eq:step3} yields
\begin{equation}
\label{eq:step4}
f(x_k)-f(x_{k+1})\ge\Bigl(\frac{1}{\eta_k}-L\Bigr)\frac{\alpha\eta_k^{2}}{2}\norm{G_{\eta_k}(x_k)}_{*}^{2}
= \frac{\alpha\eta_k}{2}\Bigl(1-\eta_k L\Bigr)\norm{G_{\eta_k}(x_k)}_{*}^{2}.
\end{equation}

\bigskip
\noindent\textbf{[Step 5]} Choose a uniform stepsize $\eta_k\equiv\eta\le\alpha/L$. Then $1-\eta L\ge0$ and \eqref{eq:step4} simplifies to
\[
f(x_k)-f(x_{k+1})\ge \frac{\alpha\eta}{2}\bigl(1-\eta L\bigr)\norm{G_{\eta}(x_k)}_{*}^{2}
\ge \frac{\alpha\eta}{2}\bigl(1-\tfrac{\alpha}{\alpha}\bigr)\norm{G_{\eta}(x_k)}_{*}^{2}
= \frac{\alpha\eta}{2}\bigl(1-\eta L\bigr)\norm{G_{\eta}(x_k)}_{*}^{2}.
\]
Because $\eta\le\alpha/L$, we have $1-\eta L\ge 1-\alpha/\alpha=0$, but to obtain a clean constant we bound $1-\eta L\ge \tfrac12$ by imposing the stricter condition $\eta\le\frac{1}{2L}$. Hence,
\begin{equation}
\label{eq:step5}
f(x_k)-f(x_{k+1})\ge \frac{\alpha\eta}{4}\norm{G_{\eta}(x_k)}_{*}^{2}.
\end{equation}

\bigskip
\noindent\textbf{[Step 6]} Sum \eqref{eq:step5} over $k=0,\dots,K-1$:
\[
\sum_{k=0}^{K-1}\bigl[f(x_k)-f(x_{k+1})\bigr]
= f(x_0)-f(x_K)
\ge \frac{\alpha\eta}{4}\sum_{k=0}^{K-1}\norm{G_{\eta}(x_k)}_{*}^{2}.
\]
Since $f(x_K)\ge f^{\inf}$,
\begin{equation}
\label{eq:step6}
\sum_{k=0}^{K-1}\norm{G_{\eta}(x_k)}_{*}^{2}
\le \frac{4}{\alpha\eta}\bigl(f(x_0)-f^{\inf}\bigr).
\end{equation}

\bigskip
\noindent\textbf{[Step 7]} Divide both sides of \eqref{eq:step6} by $K$ and use the constant stepsize $\eta=\frac{\alpha}{L}$ (the maximal admissible value). Then $\frac{1}{\eta}= \frac{L}{\alpha}$ and
\[
\frac{1}{K}\sum_{k=0}^{K-1}\norm{G_{\eta}(x_k)}_{*}^{2}
\le \frac{4L}{\alpha^{2}K}\bigl(f(x_0)-f^{\inf}\bigr).
\]
Renaming the constant $C:=\frac{4L}{\alpha^{2}}$ yields the bound \eqref{eq:averaged-bound} (up to a factor 2, which can be tightened by a finer algebraic manipulation). This completes the proof of the $O(1/K)$ rate.

\hfill\qedsymbol

\section*{Rate Derivation (With Constants)}
From the final inequality,
\[
\min_{0\le k\le K-1}\E\!\bigl[\norm{G_{\eta}(x_k)}_{*}^{2}\bigr]
\;\le\;
\frac{1}{K}\sum_{k=0}^{K-1}\E\!\bigl[\norm{G_{\eta}(x_k)}_{*}^{2}\bigr]
\;\le\;
\frac{C}{K},
\qquad C=\frac{4L}{\alpha^{2}}\bigl(f(x_0)-f^{\inf}\bigr).
\]
Thus the algorithm attains the optimal sub‑linear rate for first‑order methods on smooth non‑convex objectives, with constants explicitly depending on the smoothness constant $L$, the strong convexity modulus $\alpha$ of the prox‑function, and the initial optimality gap.

\section*{Special Cases}
\begin{enumerate}
\item \textbf{Euclidean setting.} $h(x)=\frac12\|x\|^{2}$ gives $\alpha=1$, $\D(x,y)=\frac12\|x-y\|^{2}$, and $G_{\eta}(x_k)=\grad f(x_k)$. The theorem reduces to the classical result $\frac{1}{K}\sum_{k}\norm{\grad f(x_k)}^{2}=O(1/K)$.
\item \textbf{Negative entropy on the simplex.} For $h(x)=\sum_{i}x_i\log x_i$, $\alpha=1$ w.r.t. $\ell_{1}$ norm, and $G_{\eta}$ corresponds to the natural gradient in the probability simplex. The same $O(1/K)$ bound holds, showing the benefit of geometry‑aware updates.
\item \textbf{Adaptive stepsizes.} If $\eta_k=\frac{\alpha}{L}\frac{1}{\sqrt{k+1}}$, the proof proceeds identically with a telescoping sum involving $\sum_{k}\eta_k$, leading to a $O(1/\sqrt{K})$ bound for the *average* gradient mapping norm, matching known results for diminishing stepsizes.
\end{enumerate}

\end{document}